\documentclass[11pt]{article}

\usepackage[margin=1in]{geometry}
\usepackage{amsmath, amssymb}
\usepackage{booktabs}
\usepackage{graphicx}
\usepackage{hyperref}
\usepackage{url}
\usepackage{enumitem}
\setlist{leftmargin=*}
\setlength{\parskip}{0.35em}
\setlength{\parindent}{0pt}

\title{\textbf{Stage 4: Experimental Evaluation}\\
Runtime, Scalability, and Baseline Comparison}
\author{Team Members: Name 1, Name 2, Name 3}
\date{CS240: Algorithm Design and Analysis, Spring 2026}

\begin{document}

\maketitle

\section{Evaluation Goals}

The goal of this stage is to evaluate whether the implemented algorithms
actually support the project claim: classical graph partitioning algorithms can
reduce communication cost in distributed optimization while keeping workloads
balanced. The evaluation focuses on three requirements from the course
announcement:
\begin{enumerate}
    \item measuring and analyzing runtime;
    \item evaluating scalability as graph size increases;
    \item comparing the reproduced method against at least one baseline or
    variant.
\end{enumerate}

The reproduced method is Kernighan--Lin refinement. The two baselines are
greedy balanced partitioning and recursive spectral bisection. The greedy
method is a fast constructive baseline, while spectral bisection is a stronger
algorithmic baseline based on graph Laplacian eigenvectors.

\section{Experimental Setup}

All experiments are run through the shared experiment driver:
\[
\texttt{python src/experiments.py --output-dir outputs --seed 7}.
\]
The output is stored in \texttt{outputs/partition\_results.csv}. The same graph
instances, partition counts, metrics, and random seed are used for all methods.
This makes the comparison direct and reproducible.

The benchmark suite contains:
\begin{itemize}
    \item grid graphs with side lengths 6, 8, and 10;
    \item Erdos--Renyi graphs with 40 and 80 vertices;
    \item random geometric graphs with 50 and 90 vertices;
    \item the Zachary karate club graph.
\end{itemize}

For most synthetic graphs, we use $k=4$ partitions. For the karate graph, we
use $k=2$ because the graph is small and naturally associated with a two-way
community split.

\section{Metrics}

The main partition-quality metric is cut weight:
\[
    \mathrm{cut}(p)
    =
    \sum_{(u,v)\in E,\;p(u)\neq p(v)} w_{uv}.
\]
Lower cut weight means fewer cross-worker dependencies.

The balance metric is
\[
    \rho(p)=
    \frac{\max_i |\{v:p(v)=i\}|}{|V|/k}.
\]
A value of $1$ means perfect balance. Values close to $1$ indicate that the
algorithm avoids assigning too many vertices to one worker.

Runtime is measured in seconds using the wall-clock time required to compute
the partition. In addition, a consensus simulation reports cross-partition
message counts. Since each edge crossing a worker boundary creates two
directed messages per iteration in the simulation, lower cut weight should
produce lower cross-worker traffic.

\section{Cut-Weight Comparison}

Across all benchmark instances, the average cut weights are:
\[
    \text{greedy}=98.25,\qquad
    \text{spectral}=72.75,\qquad
    \text{Kernighan--Lin}=66.88.
\]
Thus, Kernighan--Lin refinement has the best average partition quality. It
improves substantially over the greedy baseline and also improves slightly
over the spectral baseline on average.

\begin{center}
\begin{tabular}{llrrr}
\toprule
Graph & Size & Greedy & Spectral & Kernighan--Lin \\
\midrule
grid & 6 & 23 & 19 & 18 \\
grid & 8 & 40 & 20 & 20 \\
grid & 10 & 65 & 32 & 29 \\
erdos\_renyi & 40 & 54 & 44 & 37 \\
erdos\_renyi & 80 & 122 & 151 & 119 \\
geometric & 50 & 179 & 136 & 136 \\
geometric & 90 & 278 & 155 & 153 \\
karate & 34 & 25 & 25 & 23 \\
\bottomrule
\end{tabular}
\end{center}

The results show several patterns. On grid graphs, spectral bisection is
already strong because the graph has clear geometric structure, but
Kernighan--Lin can still improve or match it. On random geometric graphs, both
spectral and Kernighan--Lin perform much better than greedy partitioning. On
the larger Erdos--Renyi graph, spectral partitioning performs worse than the
greedy baseline, showing that spectral structure is not always aligned with
low balanced cut in random graphs; Kernighan--Lin refinement recovers a better
solution.

\section{Runtime Analysis}

The greedy baseline is consistently the fastest method. This is expected
because it constructs a partition in one pass and does not solve eigenvalue
problems or perform local-search refinement.

Spectral bisection is slower than greedy because it computes eigenvectors of
the graph Laplacian. In the current implementation, dense eigen-decomposition
is used, so the cost grows quickly with graph size. For the course-scale
graphs, this is acceptable; for much larger graphs, sparse eigensolvers would
be necessary.

Kernighan--Lin refinement is usually the slowest because it starts from a
spectral partition and then performs additional pairwise refinement. The extra
runtime is the cost paid for lower cut values. This is a reasonable tradeoff
when communication cost is more important than partitioning time, especially if
partitioning is performed once and the resulting worker assignment is reused
over many distributed optimization iterations.

\section{Scalability Evaluation}

Scalability can be observed by comparing instances within the same graph
family. For grid graphs, the number of vertices increases from 36 to 64 to
100, and the number of edges increases from 60 to 112 to 180. The cut weights
increase with graph size for all methods, but the relative ordering remains
stable: greedy has the largest cut, spectral is much lower, and
Kernighan--Lin is the best or tied for best.

For Erdos--Renyi graphs, increasing from 40 to 80 vertices makes the graph
denser and more difficult. The larger instance shows that no single baseline
is always reliable: spectral bisection has cut weight 151, worse than greedy's
122, while Kernighan--Lin obtains 119. This suggests that local refinement is
valuable as a corrective step when the initial spectral direction is not ideal.

For random geometric graphs, increasing from 50 to 90 vertices also increases
edge count substantially. Spectral partitioning remains strong because spatial
structure is meaningful, and Kernighan--Lin slightly improves the larger
instance from 155 to 153. This indicates that local refinement may give modest
but consistent improvements on graphs whose global geometry is already well
captured by spectral methods.

\section{Communication Evaluation}

The consensus simulation reports cross-partition message counts. The average
cross-worker messages are:
\[
    \text{greedy}=15460,\qquad
    \text{spectral}=11360,\qquad
    \text{Kernighan--Lin}=10440.
\]
These numbers follow the same trend as cut weight. This is expected because
cross-partition messages are generated by edges whose endpoints lie in
different parts. Lower cut weight therefore corresponds to lower communication
placement cost.

The final consensus error is nearly identical across different partitions on
the same graph. This happens because the simulation keeps the original graph
topology fixed and uses the partition only to count whether messages cross
worker boundaries. This design isolates communication cost from convergence
dynamics. A more advanced simulation could make cross-worker communication
slower, delayed, or budget-limited so that partitions also affect convergence.

\section{Figures}

The generated figures summarize the main experimental results:
\begin{figure}[ht]
    \centering
    \includegraphics[width=0.78\linewidth]{outputs/figures/cut_weight.png}
    \caption{Cut-weight comparison. Lower values indicate less cross-worker
    communication.}
\end{figure}

\begin{figure}[ht]
    \centering
    \includegraphics[width=0.78\linewidth]{outputs/figures/runtime_seconds.png}
    \caption{Runtime comparison. Greedy is fastest; Kernighan--Lin spends
    additional time on refinement.}
\end{figure}

\begin{figure}[ht]
    \centering
    \includegraphics[width=0.78\linewidth]{outputs/figures/balance_ratio.png}
    \caption{Balance-ratio comparison. All methods maintain good balance in
    the benchmark suite.}
\end{figure}

\section{Conclusion}

The experimental results satisfy the evaluation requirements. Runtime is
measured, scalability is examined through increasing graph sizes, and
Kernighan--Lin refinement is compared against two baselines. The results show
that the reproduced refinement method usually provides the lowest cut and the
lowest cross-worker communication, at the cost of additional runtime. This
supports the project's main claim that classical graph partitioning algorithms
can help organize distributed computation more communication-efficiently.

\end{document}
